\section{Ontologies and Knowledge Representation}
\label{sec:ontologies-and-knowledge-representation}

This section presents the background on ontologies and knowledge
representation that the system in this thesis builds on. It first
describes what an ontology is and how one is structured, then
explains why the ML domain is suited to ontological
representation. It then introduces the Semantic Web standards used
to implement ontologies, and reviews existing ML ontologies. The
last part examines a specific ontology, called ML Ontology, in more detail, since it is used as
a starting point for the work in this thesis.

\subsection{Ontology Description}
\label{sec:ontology-description}

\textcite[1]{gruber1995toward_principles_design_ontologies_used_knowledge_sharing}
define an ontology as
\enquote{an explicit specification of a conceptualization},
where a conceptualisation is a simplified abstraction of the part of the
world being modelled, describing which entities are taken to exist in it
and how they are related.
\textcite[184]{studer1998knowledge_engineering_principles_methods} refine
this, describing an ontology as a
\enquote{formal, explicit specification of a shared conceptualisation}.
Each term in this definition carries a distinct meaning.
\textit{Formal} means that the ontology is machine-readable rather than
expressed in natural language.
\textit{Explicit} means that the concepts and the constraints on how
they may be used are defined directly rather than left implicit.
\textit{Shared} means that the ontology captures knowledge accepted by a
group rather than knowledge private to a single individual
\parencite{studer1998knowledge_engineering_principles_methods}.
In practical terms, an ontology gives a domain a shared set of terms,
together with definitions of its basic concepts and their relations
\parencite{noy2001ontology_development_101_guide_creating_first_ontology}.

An ontology consists of classes, which represent concepts in the
domain. It also has properties, or slots, that describe the features
of those concepts, and restrictions on those properties that
constrain their allowed values
\parencite{noy2001ontology_development_101_guide_creating_first_ontology}.
Classes are typically organised in a hierarchy where subclasses
inherit the properties of their parent classes. Taken together, these
elements allow a domain to be described in a way that both humans and
software can interpret and reason over
\parencite{noy2001ontology_development_101_guide_creating_first_ontology}.

The design of an ontology involves several practical constraints.
\textcite{gruber1995toward_principles_design_ontologies_used_knowledge_sharing}
argues that an ontology should make the minimal ontological commitment
needed for its purpose, committing to as little
about the modelled world as the task allows. This connects to a
well-known trade-off in ontology design. The more reusable an
ontology is, the less directly usable it tends to be for a specific
task, because a more general representation requires more refinement
before it can be applied
\parencite{studer1998knowledge_engineering_principles_methods}.
\textcite{noy2001ontology_development_101_guide_creating_first_ontology}
point to a related constraint on scope. An ontology should not attempt
to capture everything that could be said about a domain, but only what
is needed for the application at hand
\parencite{noy2001ontology_development_101_guide_creating_first_ontology}.
Expanding scope beyond what the
application requires adds complexity without benefit.
Ontology development is also not a one-time activity. As domain
knowledge evolves and new requirements emerge, the ontology must be
revised, which means the initial design should anticipate extension
and allow the vocabulary to be specialised without requiring changes
to existing definitions
\parencite{gruber1995toward_principles_design_ontologies_used_knowledge_sharing,
noy2001ontology_development_101_guide_creating_first_ontology}.

\subsection{Machine Learning and Ontology}
\label{sec:ml-and-ontology}

The machine learning domain has several properties that make 
ontological representation a natural fit.

ML methods are not a flat collection of independent techniques.
They form hierarchical structures where more specific methods are 
specialisations of more general ones.
\textcite{humm2026industry_ready_machine_learning_ontology} 
illustrates this directly: a multilayer perceptron is a 
specialisation of an artificial neural network, and if the latter 
can be used for classification and regression, the former inherits 
that applicability without needing it to be stated separately.
This kind of reasoning over hierarchical structure is something an 
ontology supports naturally through inheritance 
\parencite{noy2001ontology_development_101_guide_creating_first_ontology}.

Beyond hierarchical structure, ML concepts are connected through 
typed semantic relations.
Methods are suited to certain tasks and data types, belong to certain 
learning paradigms, and have different performance characteristics 
\parencite{humm2026industry_ready_machine_learning_ontology}.
Making these associations explicit allows methods to be compared
in a structured way. 
\textcite{keet2015data_mining_optimization_ontology} motivate their
ontology for data mining on these grounds. In traditional
meta-learning, learning algorithms were treated as black boxes, and
only the characteristics of the data were linked to observed
performance. An ontology instead allows the internal workings of an
algorithm to be described explicitly.

\subsection{Semantic Web Standards}
\label{sec:semantic-web-standards}

Ontologies, both in the ML domain and in general, are commonly implemented using
standardised languages from the Semantic Web. This is a set of
specifications developed by the World Wide Web Consortium
(W3C) to make information machine-readable and interoperable
across systems \parencite{w3c}.

The foundation is the Resource Description Framework (RDF),
the standard knowledge representation language for the Semantic
Web \parencite{gibbins2009resource_description_framework_rdf}.
RDF represents knowledge as triples, each consisting of a
subject, a predicate, and an object, forming a directed graph
of binary relationships
\parencite{gibbins2009resource_description_framework_rdf}.
RDF Schema (RDFS) extends RDF with basic constructs for
defining classes, properties, and constraints
\parencite{gibbins2009resource_description_framework_rdf}.
Building on RDF, the Simple Knowledge Organization System (SKOS)
provides a lightweight standard for representing taxonomies
and controlled vocabularies
\parencite{humm2026industry_ready_machine_learning_ontology}.

The Web Ontology Language (OWL) builds on RDF and RDFS and
adds richer modelling constructs, including logical operators,
quantifiers, and property characteristics such as transitivity
\parencite{bechhofer2009owl_web_ontology_language}.
This gives OWL well-defined semantics that support automated
reasoning
\parencite{bechhofer2009owl_web_ontology_language}.

The SPARQL Protocol and RDF Query Language (SPARQL) is used to
query RDF data. It has syntax similar to Structured Query Language (SQL) and retrieves
data by matching triple patterns against stored graphs
\parencite{grobe2009rdf_jena_sparql_semantic_web}.
RDF data is typically stored in triplestores, database systems
designed specifically for handling RDF triples
\parencite{grobe2009rdf_jena_sparql_semantic_web}.

\subsection{Existing Machine Learning Ontologies}
\label{sec:existing-machine-learning-ontologies}

Several ontologies have been developed to represent knowledge
in the ML and data mining domain, each with a different scope
and purpose. The following is a selection of openly available
ontologies identified within the ML domain.

ML-Schema \parencite{publio2018ml_schema_exposing_semantics_machine_learning_schemas_ontologies},
developed by the W3C Machine Learning Schema Community Group, is a
top-level schema that defines shared classes and properties for
describing and sharing metadata on ML algorithms, datasets, models,
and experiments.
Its primary goal is interoperability. It provides a shared
reference vocabulary that other ML ontologies can be mapped to,
reducing fragmentation across platforms.
ML-Schema defines around 25 abstract classes but contains no
populated instances, meaning it describes structure rather than
ML content.
It is intended as a foundation for more specific ontologies,
not as a standalone knowledge base.

OntoDM-core \parencite{panov2014ontology_core_data_mining_entities}
is a formal ontology for core data mining entities such as
datasets, tasks, algorithms, and their executions.
It is structured across three layers. A specification layer
covers abstract definitions, an implementation layer covers
algorithm implementations and parameters, and an application
layer covers executions and workflows.
OntoDM-core is aligned with established top-level ontologies,
which enables its use across different application domains
\parencite{panov2014ontology_core_data_mining_entities}.
With 856 classes and detailed taxonomies of tasks and
algorithms, it provides high structural granularity 
for data mining concepts.

The MEX Vocabulary \parencite{esteves2015mexvoc} is a
lightweight format for exchanging provenance metadata about
ML experiments.
It is organised in three layers. MEX-Core handles experiment
configuration and execution, MEX-Algorithm covers algorithm
class and learning method, and MEX-Performance covers
evaluation measures.
MEX is designed for basic ML iterations only and does not aim
to cover the full ML domain.

ML Ontology \parencite{humm2026industry_ready_machine_learning_ontology}
is an ontology for representing ML content. It covers ML areas,
tasks, approaches, preprocessing methods, metrics, libraries,
Automated Machine Learning (AutoML) solutions, and configuration parameters, with around 700
instances and 5000 RDF triples. ML Ontology
is designed for integrating
different ML tools and frameworks through a shared vocabulary. It
supports forward-chaining inference via SPARQL and
includes built-in quality assurance checks.

It is clear that these ontologies have different aims. ML-Schema and
MEX are built for annotating and exchanging experiment metadata.
OntoDM-core is a broader reference ontology for data mining
entities such as datasets, tasks, and algorithms. OntoDM-core can be
queried to find algorithms that suit a given task and data type, but
its coverage is limited to the data mining domain. ML Ontology covers
a wider range of ML areas, methods, tasks, and metrics, and is
populated accordingly, which makes it a more suitable basis for
reasoning about method selection across the ML domain.

\subsection{ML Ontology Structure and Content}
\label{sec:ml-ontology-structure-and-content}

Among the ontologies reviewed, ML Ontology by \textcite{humm2026industry_ready_machine_learning_ontology}
is examined in
more detail here because it provides the broadest coverage
of ML concepts and methods, making it the most promising
foundation for the system developed in this thesis.

ML Ontology is organised into five modules: classes,
instances, inference rules, quality assurance checks, and
predefined queries.
Figure~\ref{fig:ml-ontology-schema} shows this overall
structure.

\begin{figure}[H]
    \centering
    \includegraphics[width=0.8\textwidth]{Figures/02-background/ml-ontology-schema.png}
    \caption[ML Ontology schema]{ML Ontology schema. Reproduced from
    \textcite{humm2026industry_ready_machine_learning_ontology},
    licensed under CC BY 4.0.}
    \label{fig:ml-ontology-schema}
\end{figure}

A few naming conventions in ML Ontology are worth noting.
The class called \texttt{ML\_area} represents the same concept as the learning paradigm 
described in Section~\ref{sec:machine-learning-paradigms-and-tasks}, with
instances for \textit{supervised\_learning}, \textit{unsupervised\_learning}, and \textit{reinforcement\_learning}.
The \texttt{ML\_approach} class represents ML methods, with instances for specific algorithms such as \textit{random\_forest} and \textit{support\_vector\_machine}.

The classes are divided into two groups: ML concepts, which
cover the domain vocabulary, and ML implementations, which
cover libraries, AutoML solutions, and their configuration
options.
Figure~\ref{fig:ml-ontology-classes} shows the main classes
and the relations between them.

\begin{figure}[H]
    \centering
    \includegraphics[width=\textwidth]{Figures/02-background/ml-ontology-classes.png}
    \caption[ML Ontology classes]{ML Ontology classes. Reproduced from
    \textcite{humm2026industry_ready_machine_learning_ontology},
    licensed under CC BY 4.0.}
    \label{fig:ml-ontology-classes}
\end{figure}

Some key relations are presented in Figure~\ref{fig:ml-ontology-classes}, 
among them are \texttt{used\_for}, which links
approaches to the tasks they apply to.
Tasks are also linked to ML areas via \texttt{belongs\_to}.
\texttt{can\_perform} links libraries and AutoML
solutions to tasks. Instances are also linked to corresponding 
entries in Wikidata \parencite{wikidata} via \texttt{rdfs:seeAlso}.

Since ML methods are central to the method selection task in this 
thesis, the properties of \texttt{ML\_approach} instances are 
examined in more detail. Table~\ref{tab:ml-approach-properties} 
lists all properties associated with \texttt{ML\_approach} instances 
in ML Ontology, as found in the published Turtle file 
\parencite[Section~6.7]{humm2026industry_ready_machine_learning_ontology}.

\begin{table}[H]
\centering
\caption[ML\_approach properties]{Properties of \texttt{ML\_approach} 
instances in ML Ontology.}
\label{tab:ml-approach-properties}
\begin{tabularx}{\textwidth}{lX}
\toprule
\textbf{Property} & \textbf{Description} \\
\midrule
\texttt{skos:prefLabel} & The preferred display name of the approach. \\
\midrule
\texttt{skos:altLabel} & Alternative names or abbreviations, 
for example \textit{ANN} for artificial neural network. \\
\midrule
\texttt{skos:broader} & Links the approach to a more general 
parent approach in the hierarchy. \\
\midrule
\texttt{rdfs:comment} & A short human-readable description 
of the approach. \\
\midrule
\texttt{rdfs:seeAlso} & A link to an external resource, 
typically a Wikidata entry. \\
\midrule
\texttt{used\_for} & The ML tasks this approach can be 
applied to. \\
\midrule
\texttt{performance} & Performance characteristics, 
for example \textit{fast\_training} or \textit{accurate}. \\
\midrule
\texttt{possible\_if} & Conditions under which the approach 
is applicable, for example \textit{numerical\_features}. \\
\midrule
\texttt{not\_recommended\_if} & Conditions under which the 
approach should be avoided, for example \textit{many\_features}. \\
\bottomrule
\end{tabularx}
\end{table}

Coverage varies considerably across properties.
Table~\ref{tab:ml-approach-coverage} shows how many of the 122
\texttt{ML\_approach} instances have each property populated, based
on an analysis of the published Turtle file
\parencite[Section~6.7]{humm2026industry_ready_machine_learning_ontology}.
This coverage of the original ontology serves as a baseline for the
extensions made in this thesis.

\begin{table}[H]
\centering
\caption[ML\_approach property coverage]{Coverage of 
\texttt{ML\_approach} properties across 122 instances in 
ML Ontology.}
\label{tab:ml-approach-coverage}
\begin{tabular}{lr}
\toprule
\textbf{Property} & \textbf{Coverage} \\
\midrule
\texttt{skos:prefLabel}       & 122/122 (100\%) \\
\texttt{used\_for}           & 103/122 (84\%)  \\
\texttt{rdfs:comment}         & 97/122 (79\%)   \\
\texttt{skos:altLabel}        & 64/122 (52\%)   \\
\texttt{skos:broader}         & 56/122 (45\%)   \\
\texttt{rdfs:seeAlso}         & 56/122 (45\%)   \\
\texttt{performance}         & 23/122 (18\%)   \\
\texttt{possible\_if}        & 9/122 (7\%)     \\
\texttt{not\_recommended\_if}& 3/122 (2\%)     \\
\bottomrule
\end{tabular}
\end{table}

The low coverage of some properties reflects that they must 
be populated explicitly for each instance. However, ML Ontology 
includes a forward-chaining inference rule that propagates 
selected properties along \texttt{skos:broader} hierarchies. 
Properties marked as \textit{broader inherited} are 
automatically assigned to more specific concepts based on 
what is defined for their parent. For the properties in 
Table~\ref{tab:ml-approach-coverage}, this applies to 
\texttt{used\_for},
\texttt{performance}, \texttt{possible\_if}, and 
\texttt{not\_recommended\_if}. The rule must be executed 
after loading the ontology, and is not applied automatically.

ML tasks are also organised in hierarchies using 
\texttt{skos:broader}. For tasks, the inference rule means 
that if \texttt{classification} belongs to 
\texttt{supervised\_learning}, then \texttt{text\_classification} 
inherits that relation without it needing to be stated 
explicitly. Each \texttt{ML\_task} instance may also carry a 
\texttt{has\_dataset\_type} property that specifies which data 
types the task applies to, for example \textit{tabular} or 
\textit{image}. Since \texttt{ML\_approach} instances are linked 
to tasks via \texttt{used\_for}, dataset type is accessible 
indirectly through this relation. This means that an approach 
applicable to \texttt{image\_classification} is implicitly 
associated with image data through the task it is linked to. 
Figure~\ref{fig:ml-ontology-task-hierarchy} illustrates 
the task hierarchy for some tasks that belong to 
supervised and unsupervised learning. 

\begin{figure}[H]
    \centering
    \includegraphics[width=\textwidth]{Figures/02-background/ml-ontology-task-hierarchy.png}
    \caption[ML task hierarchy]{ML task hierarchy in ML Ontology. Reproduced from
    \textcite{humm2026industry_ready_machine_learning_ontology},
    licensed under CC BY 4.0.}
    \label{fig:ml-ontology-task-hierarchy}
\end{figure}

Figure~\ref{fig:ml-ontology-task-hierarchy} shows the task hierarchy for a set
of classification, regression, clustering and feature reduction tasks. 
We can see that many of the tasks are narrower versions of more general ones, 
and differ in the data type they apply to. 
For example, \texttt{tabular\_classification} is a specialisation of \texttt{classification}.

There are some missing concepts in the ontology. 
Reinforcement learning is defined as an ML area in the 
ontology but does not appear in the hierarchy in Figure~\ref{fig:ml-ontology-task-hierarchy}
since no tasks
and no methods have been assigned to it.
Coverage of \texttt{ML\_approach} instances is also
incomplete. \textcite{humm2026industry_ready_machine_learning_ontology}
report a recall of 0.84 for this class, with graph neural
networks, variational autoencoders, and genetic algorithms
identified as examples of missing concepts.